
%%%%%%%%%%%%%%%%%%%%%%%%%%%%%%%%%%%%%%%%%%%%%%%%%%%%%%%
\begin{frame}{Tri externe}

Le \alert{tri externe} est utilisé,

\begin{redbullet}
  \item pour les algorithmes de jointure (\textit{sort/merge})

   \item l'élimination des doublons (clause \texttt{distinct})

    \item pour les opérations de regroupement (\texttt{group by})

   \item ... et bien sûr pour les \texttt{order by}
\end{redbullet}

C'est une opération qui peut être très coûteuse
sur de grands jeux de données.
\end{frame}
%%%%%%%%%%%%%%%%%%%%%%%%%%%%%%%%%%%%%%%%%%%%%%%%%%%%%%%
\begin{frame}{Première phase\,: le tri}

On remplit la mémoire, on trie, on vide
dans des \alert{fragments}, et
on recommence.

\medskip

\figSlide{phasetri}

\medskip

Coût\,: une lecture + une écriture du fichier.

\end{frame}
%%%%%%%%%%%%%%%%%%%%%%%%%%%%%%%%%%%%%%%%%%%%%%%%%%%%%%%
\begin{frame}{Fusion de listes triées}

\figSlideWithSize{fusion}{9cm}

\end{frame}

\begin{frame}{Deuxième phase\,: la fusion}

On groupe les fragments par $M$ (taille de la zone
de tri), et on fusionne.

\medskip


\figSlideWithSize{trifusion}{9cm}

\medskip

Coût\,: autant de lectures/écritures du fichier que de niveaux
de fusion.

\end{frame}

%%%%%%%%%%%%%%%%%%%%%%%%%%%%%%%%%%%%%%%%%%%%%%%%%%%%%%%
\begin{frame}{Illustration avec $M = 3$}

\figSlideWithSize{tri-fusion}{9cm}

\end{frame}
%%%%%%%%%%%%%%%%%%%%%%%%%%%%%%%%%%%%%%%%%%%%%%%%%%%%%%%
\begin{frame}{Essentiel\,: la taille de la zone de tri}

Un fichier de 75~000 pages de 4 Ko, soit 307 Mo.
\begin{redbullet}
  \item \textbf{$M > 307 Mo$}\,: une lecture, soit 307 

  \item \textbf{$M = 2 Mo$}, soit 500 pages.
    \begin{itemize}
      \item le tri donne $\lceil \frac{307}{2} \rceil = 154$
        fragments.

      \item On fait la fusion avec $154$ pages
     \end{itemize}
       Coût total de $614 + 307 = 921$ Mo. 
\end{redbullet}

NB\,: il faut allouer beaucoup de mémoire pour
passer de 1 à 0 niveau de tri.

\end{frame}


\begin{frame}{Avec très peu de mémoire}

\textbf{$M = 1 Mo$}, soit 250 pages.

\vfill

    \begin{redbullet}
       \item on obtient 307 fragments.

       \item On fusionne les 249 premiers fragments,
               puis les 58 restant. On obtient $F_1$ et $F_2$.
        \item On fusionne $F_1$ et $F_2$.
   \end{redbullet}

\vfill

 Coût total\,: $1\,228 + 307 = 1\,535$ Mo.
\vfill
Résultat\,: grosse dégradation entre 2 Mo et 1 Mo (calcul
approximatif).



\end{frame}


%%%%%%%%%%%%%%%%%%%%%%%%%%%%%%%%%%%%%%%%%%%%%%%%%%%%%%%
\begin{frame}[fragile]{L'opérateur de tri}

L'opérateur de tri est \alert{bloquant}

\begin{redbullet}
  \item Les deux phases du tri-fusion sont effectuées pendant le \textit{open()}
  \item Le \textit{next()} ne fait que lire, un à un, les tuples dans le
     résultat du tri.
\end{redbullet}

\vfill

Conséquence: \alert{latence importante des requêtes impliquant un tri}

\vfill

Grosse différence entre:

\begin{minted}[frame=leftline,
               baselinestretch=.9,
               fontsize=\footnotesize,
               framesep=2mm]{sql}
select * from Film
\end{minted}
et
\begin{minted}[frame=leftline,
               baselinestretch=.9,
               fontsize=\footnotesize,
               framesep=2mm]{sql}
select * from Film order by titre
\end{minted}
\end{frame}

